\documentclass[11pt,letterpaper]{article}
\usepackage[margin=1in]{geometry}
\usepackage[T1]{fontenc}
\usepackage{lmodern}
\usepackage{microtype}
\usepackage{amsmath,amssymb}
\usepackage{booktabs,array,enumitem,xcolor}
\usepackage{titlesec}
\usepackage[colorlinks=true,linkcolor=black,urlcolor=blue!50!black,citecolor=black]{hyperref}
\setlist{itemsep=2pt,topsep=3pt}
\setlength{\parindent}{0pt}\setlength{\parskip}{6pt}
\titleformat{\section}{\normalfont\large\bfseries}{\thesection}{0.6em}{}
\titlespacing*{\section}{0pt}{12pt}{4pt}
\pagestyle{plain}
\newcommand{\plainwords}[1]{\par\medskip\noindent\fcolorbox{black!25}{black!4}{\parbox{\dimexpr\linewidth-2\fboxsep-2\fboxrule}{\textbf{In plain words.} #1}}\par\medskip}
\newcommand{\titleblock}[3]{%
  {\LARGE\bfseries #1\par}\vspace{4pt}
  {\large #2\par}\vspace{2pt}
  {\small\color{black!70} Shakil A.\ Rafi, Ph.D.\ \textbar\ University of Arkansas \textbar\ #3\par}\vspace{2pt}
  \hrule\vspace{6pt}}

\begin{document}
\titleblock{Semantic Interrogation Encoding}{How many questions does a document need?}{Research write-up, November 2026}
\textit{In collaboration with Johns Hopkins University. Status: active research. Code is private and available to reviewers on request.}

\plainwords{Instead of turning a text into a block of anonymous numbers, I ask a language model a list of questions about it and keep its answers. Every number then has a meaning I can read. What I am testing is how short that list can get before it stops being useful.}

\section{Introduction}
A dense embedding is compact and works well, but its coordinates carry no declared meaning. For a text $x$, a question bank $Q$, and a probabilistic decision model $M$, the \emph{Semantic Interrogation Encoding} (SIE) is the concatenation of the decision distributions $M(x,q_i)$. Each coordinate is the model's answer to a named question.

The claim is deliberately modest: a compact, interpretable question bottleneck may preserve useful predictive information. I am not claiming that question answering beats a dense embedding. Prior work overlaps substantially with the embedding idea, so what I am after is careful measurement of the accuracy and cost trade-off, not novelty of the concept.

\section{Method}
\textbf{Decision model.} Questions are answered by Laya (\texttt{laya-mlx} 0.3.0), run offline. Questions are typed (binary, categorical, ordinal); binary questions contribute one probability, the others keep their full distributions.

\textbf{Question banks.} I use two. The first is a deterministic 128-question template bank across 32 semantic families, oriented toward sentiment. The second is learned from unlabeled text: I sample FineWeb-Edu, embed it with a pinned MiniLM revision, cluster it into 192 candidates, keep the 128 best-supported clusters, and write one topic question per cluster from cluster-distinctive terms. No downstream labels are used to build either bank.

\textbf{Classifier and controls.} One small logistic regression is used identically for every representation. I report nested curves at 8, 16, 32, 64, and 128 questions, and compare against five random question subsets, a deliberately meaningless question bank, a signed-hashing baseline, MiniLM, and label-free random projections of matching dimension. Intervals are 95\% paired bootstraps, and I report question count, feature dimension, tokens, and latency.

\textbf{Rules I hold myself to.} Freeze the split before encoding; build questions without evaluation labels; reject input truncation rather than hide it (Laya has a 512-token limit); fit selection and classifiers on training data only.

\section{Results}
\begin{center}
\small
\begin{tabular}{@{}p{3.3cm}p{6.4cm}p{2.6cm}@{}}
\toprule
\textbf{Experiment} & \textbf{Result} & \textbf{Status}\\
\midrule
SST-2 & 91.4\% accuracy at 32 questions & Pushed\\
Rotten Tomatoes & 88.8\% at 64 questions, same frozen bank & Pushed\\
Sentiment cascade (SST-2) & $\approx$89\% at $\approx$12.7 questions on average; about $3\times$ faster than fixed-32 & Preliminary\\
Yelp transfer (frozen) & 94.5\% hierarchy accuracy; a fixed 8-question encoding scored 95.3\% & Preliminary\\
DynaSent Round 2 (binary) & $\approx$78.5--79\%; confident errors; no neutral class evaluated & Preliminary\\
\bottomrule
\end{tabular}
\end{center}
The first two numbers are the best results among the question counts I evaluated, so they are descriptive, not independently confirmed model selections. The cascade, Yelp, and DynaSent results are single runs from my working notes and have not been through the same review.

\section{Discussion}
Reaching 91.4\% on SST-2 with 32 questions shows that a compact, readable code can carry sentiment well. It does not show more than that. The bank is sentiment-oriented, the texts are short, and the two pilots do not demonstrate transfer between task types.

The cascade suggests that most documents need far fewer than 32 questions. The Yelp result is a useful caution, though: the plain fixed 8-question encoding matched or beat the cascade, which tells me confidence alone is not a validated stopping rule. The MiniLM baseline was weaker than I expected and needs validating before I make any comparison claim, and MiniLM is still much faster. Laya jointly encodes each text and question, so the document representation is not reusable across arbitrary questions. Any systems-cost argument has to face that directly.

\textbf{What would weaken the idea.} SIE needing comparable effective dimension while staying materially worse; optimized banks failing to beat random ones; transfer collapsing across tasks; reconstruction retaining little beyond labels; repeated interrogation losing the cost comparison; or unstable, poorly calibrated probabilities. I am keeping the negative results.

\textbf{Next.} Two branches. One tests whether a corpus-learned bank transfers across tasks. The other tests whether early stopping preserves sentiment accuracy at a measured latency. Reconstruction, a VAE, and adaptive policies remain future work.

\textbf{Contact.} sarafi@uark.edu
\end{document}
